% Histogramme de vingt mesures de la période d'un pendule ; au corrigé, moyenne et classes centrales (exercices 1 et 6, chapitre 21)
\begin{tikzpicture}
\begin{axis}[otMesures,otAxesBout,width=96mm,height=46mm,xmin=1.37,xmax=1.515,ymin=0,ymax=7.5,
  xtick={1.38,1.40,...,1.50},ytick={0,1,...,6},grid=none,
  xticklabel style={/pgf/number format/fixed,/pgf/number format/fixed zerofill,/pgf/number format/precision=2},
  xlabel={$T$ (s)},ylabel={effectif}]
\solution{\fill[otGreen!20] (axis cs:1.42,0) rectangle (axis cs:1.46,6.6);}
\addplot[ybar interval,fill=otNavy!50,draw=white,line width=0.8pt] coordinates
  {(1.38,1) (1.40,3) (1.42,6) (1.44,5) (1.46,3) (1.48,2) (1.50,0)};
\solution{
  \node[pcLegende,anchor=south,text=otGreenDark] at (axis cs:1.44,6.6) {11 mesures sur 20 : \qty{55}{\percent}};
  \draw[otRed!85!black,line width=0.8pt,dash pattern=on 3pt off 2pt] (axis cs:1.4385,0) -- (axis cs:1.4385,5.6);
  \node[pcLegende,anchor=west,text=otRed!85!black] at (axis cs:1.465,4.6) {moyenne \qty{1,44}{s}};
  \draw[pcRepere] (axis cs:1.4395,4.6) -- (axis cs:1.465,4.6);}
\end{axis}
\end{tikzpicture}
\par\smallskip\legende{\enonce{Vingt élèves ont chronométré une période du même pendule.}\solution{Les mesures se regroupent autour de \qty{1,44}{s}, avec quelques valeurs plus éloignées.}}
